% 第一章 投资概要

\section{核心观点与投资逻辑}

我们对中国功率半导体行业给予\textbf{增持}评级。行业正处于 AI 算力革命与能源革命的历史性交汇点，需求端迎来三重共振——AI 数据中心功耗指数级攀升带来最大边际增量、新能源车渗透率持续提升夯实第一大下游基本盘、光伏储能装机量高增长贡献第三增长极。与此同时，第三代半导体（SiC/GaN）在高压、高频场景加速渗透，技术迭代节奏加快，国产替代从消费级向车规级纵深推进，行业成长属性显著强于周期属性。

\vspace{0.3em}

我们的核心投资逻辑可归纳为三条主线：

\begin{enumerate}[leftmargin=2em]
\item \textbf{AI 算力驱动的增量弹性主线}：AI 服务器单机架功率半导体价值量从传统机架的约 2,000–5,000 美元提升至 Blackwell 时代的 5–8 万美元，增长超 10 倍。48V 母线架构渗透、GaN VRM 导入、高压化三大技术趋势叠加，带来 VR M、DC-DC、PFC 等环节的量价齐升。
\item \textbf{新能源车 + 风光储的能源转型主线}：车规级 IGBT 国产替代空间广阔，SiC 在 800V 高压平台加速渗透；光伏储能装机量持续高增长，拉动高压 MOS、IGBT 模块需求。
\item \textbf{第三代半导体技术迭代主线}：SiC 衬底产能释放、良率提升驱动成本下行，应用从高端车型向主流车型渗透；GaN 从消费快充向数据中心、车载电源拓展，市场空间十倍级扩容。
\end{enumerate}

\vspace{0.3em}

\begin{keyinsight}[投资核心结论]
功率半导体行业处于新一轮上行周期起点，AI + 新能源双轮驱动下，行业成长性显著强于周期性。优先配置具备 IDM 制造能力、车规级认证突破、第三代半导体技术储备的龙头公司。当前估值处于历史中位偏低区间，具备较好的安全边际。
\end{keyinsight}

\section{推荐组合与评级}

我们基于\textbf{业绩确定性、产业链卡位、估值安全边际、近期催化}四大维度，对覆盖的 10 家核心公司进行综合评分排序，构建推荐组合。

\begin{exhibitbox}[表 1-1：推荐组合与评级一览]
\centering
\small
\begin{tabularx}{\textwidth}{l c c c c X}
\toprule
\textbf{公司名称} & \textbf{代码} & \textbf{市值(亿元)} & \textbf{PE(TTM)} & \textbf{投资评级} & \textbf{核心逻辑} \\
\midrule
时代电气 & 688187 & 804 & 19.6 & \textcolor{riskgreen}{\textbf{买入}} & 估值最低的功率IDM，SiC + IGBT双轮驱动 \\
斯达半导 & 603290 & 314 & 95.0 & \textcolor{riskgreen}{\textbf{增持}} & IGBT国产替代龙头，车规级突破领先 \\
华润微 & 688396 & 1,065 & 117.0 & \textcolor{riskgreen}{\textbf{增持}} & 产品线最全功率IDM，制造平台优势显著 \\
士兰微 & 600460 & 753 & 164.0 & \textcolor{riskamber}{\textbf{中性}} & IDM模式12英寸进军，短期盈利承压 \\
扬杰科技 & 300373 & 692 & 55.0 & \textcolor{riskgreen}{\textbf{增持}} & 分立器件龙头，SiC后起之秀弹性大 \\
三安光电 & 600703 & 878 & 亏损 & \textcolor{riskamber}{\textbf{中性}} & 化合物半导体全产业链，短期亏损 \\
天岳先进 & 688234 & 724 & PS~50x & \textcolor{riskamber}{\textbf{中性}} & SiC衬底龙头，高成长高估值 \\
新洁能 & 605111 & 315 & 83.0 & \textcolor{riskgreen}{\textbf{增持}} & MOSFET设计龙头，品类持续拓展 \\
东微半导 & 688261 & 75 & 35.0 & \textcolor{riskgreen}{\textbf{增持}} & 高压MOS技术领先，估值合理 \\
晶丰明源 & 688368 & 195 & 247.0 & \textcolor{riskamber}{\textbf{中性}} & 电源管理芯片龙头，AI弹性标的 \\
\bottomrule
\end{tabularx}
\sourcenote{Wind，本研究团队测算，市值及 PE 数据截至 2026 年 6 月 20 日}
\end{exhibitbox}

排名说明：
\begin{itemize}[leftmargin=2em, topsep=0pt]
\item \textbf{业绩确定性权重 35\%}：考察客户结构、订单能见度、产能利用率
\item \textbf{产业链卡位权重 25\%}：考察技术壁垒、产品结构、下游应用分布
\item \textbf{估值安全边际权重 25\%}：考察 PE、PS、PB 与历史分位及全球对标
\item \textbf{近期催化权重 15\%}：考察新产品、新客户、新产能落地节奏
\end{itemize}

\section{关键数据速览}

\begin{exhibitbox}[表 1-2：核心数据一览]
\centering
\small
\begin{tabularx}{\textwidth}{X L{2.8cm} L{2.8cm} L{2.8cm} L{2.5cm}}
\toprule
\textbf{指标} & \textbf{2025E} & \textbf{2030E} & \textbf{CAGR} & \textbf{数据来源} \\
\midrule
\rowcolor{lightgrey}
\multicolumn{5}{l}{\textit{市场规模}} \\
全球功率半导体市场 & \$540-560亿 & \$700-750亿 & 5-6\% & Yole/Omdia \\
中国功率半导体市场 & ¥1,900-2,000亿 & ¥3,000-3,200亿 & 10-12\% & SEMI/本团队 \\
全球SiC器件市场 & \$30-35亿 & \$85-100亿 & 23-25\% & Yole \\
全球GaN功率器件市场 & \$8-9亿 & \$35-43亿 & 35-40\% & Yole/TrendForce \\
AI服务器功率半导体 & \$~45亿 & \$~98亿 & ~20\% & 本团队测算 \\
\midrule
\rowcolor{lightgrey}
\multicolumn{5}{l}{\textit{国产化率}} \\
整体国产化率 & 30-35\% & 45-50\% & — & 本团队综合判断 \\
低压MOSFET & 50-70\% & 65-80\% & — & 行业协会 \\
车规级IGBT & 10-15\% & 25-35\% & — & 本团队测算 \\
SiC器件 & 35-40\% & 50-60\% & — & 本团队测算 \\
SiC衬底 & 20-30\% & 40-50\% & — & 本团队测算 \\
高端模拟/驱动芯片 & 10-15\% & 20-25\% & — & 行业协会 \\
\midrule
\rowcolor{lightgrey}
\multicolumn{5}{l}{\textit{渗透率}} \\
48V服务器架构 & ~50\% & ~85\% & — & Omdia \\
GaN VRM渗透率 & 12-18\% & 40-50\% & — & 本团队测算 \\
新能源车SiC渗透率 & ~15\% & 35-45\% & — & Yole \\
\bottomrule
\end{tabularx}
\sourcenote{Yole, Omdia, SEMI, TrendForce，本研究团队综合测算}
\end{exhibitbox}

\section{投资时钟与风险提示}

\subsection{行业位置判断}

我们认为功率半导体行业当前处于\textbf{新一轮上行周期的初期阶段}，判断依据如下：

\begin{itemize}[leftmargin=2em]
\item \textbf{需求端}：AI 数据中心资本开支持续超预期，新能源车销量保持稳健增长，光伏储能装机量维持高增速，三大下游需求共振向上。
\item \textbf{供给端}：行业经过 2023–2024 年的库存去化，渠道库存已回归健康水平，部分中高端产品出现结构性紧缺。
\item \textbf{价格端}：中低压 MOS、二极管等成熟产品价格已企稳回升，高端 IGBT、SiC 器件价格维持相对坚挺。
\item \textbf{估值端}：板块估值经历 2024 年下半年回调后，处于历史中位偏低区间，龙头公司估值具备吸引力。
\end{itemize}

\subsection{核心风险提示}

\begin{riskbox}[核心风险提示]
\begin{enumerate}[leftmargin=2em, topsep=0pt]
\item \textbf{AI 资本开支不及预期风险}：若 AI 服务器需求增速放缓，将直接影响功率半导体增量需求。
\item \textbf{新能源车销量低于预期风险}：新能源汽车是功率半导体第一大下游，销量波动直接影响行业景气度。
\item \textbf{SiC 技术迭代不及预期风险}：若 SiC 良率提升、成本下降慢于预期，渗透率可能低于预测。
\item \textbf{行业价格竞争加剧风险}：国内厂商产能扩张较快，可能引发中低端产品价格战。
\item \textbf{国际贸易政策风险}：半导体设备、材料出口管制可能影响国内产能扩张节奏。
\item \textbf{地缘政治风险}：全球供应链格局变化可能带来不可预见的供需扰动。
\item \textbf{估值波动风险}：半导体板块估值弹性大，市场情绪变化可能带来较大波动。
\end{enumerate}
\end{riskbox}
